% =============================================================================
% Proposed Solution: overall architecture, IF stage and top-level integration
% Author: Abhishek Kumar Patel (2401010020)
% Included from docs/report/main.tex (other stage subsections follow this file)
% =============================================================================

\section{Proposed Solution}

\subsection{Overall Architecture}

The proposed processor is a 32-bit, in-order, single-issue pipeline with the classic five stages
described by Patterson and Hennessy~\cite{patterson_riscv}. Fig.~\ref{fig:datapath} shows the complete
datapath. Each stage is separated from the next by a pipeline register (IF/ID, ID/EX, EX/MEM,
MEM/WB) that is updated on the rising clock edge, so that each stage works on a different
instruction in every cycle. Control signals are generated once in the ID stage and travel down the
pipeline together with the instruction that they belong to.

\begin{figure}[!t]
  \centering
  % Path relative to docs/report/main.tex
  \includegraphics[width=\columnwidth]{../../diagrams/datapath.png}
  \caption{Datapath of the proposed five-stage pipelined RISC-V processor, including the forwarding
           unit and the hazard detection unit.}
  \label{fig:datapath}
\end{figure}

The supported instructions are listed in Table~\ref{tab:isa}. The encodings follow the RISC-V
unprivileged specification~\cite{riscv_spec_unpriv}.

\begin{table}[!t]
  \centering
  \caption{Supported RV32I Instruction Subset}
  \label{tab:isa}
  \begin{tabular}{|l|l|}
    \hline
    \textbf{Type} & \textbf{Instructions} \\ \hline
    R-type        & ADD, SUB, AND, OR, XOR, SLL, SRL, SRA, SLT, SLTU \\ \hline
    I-type ALU    & ADDI, ANDI, ORI, XORI, SLTI, SLTIU, SLLI, SRLI, SRAI \\ \hline
    Load / Store  & LW, SW \\ \hline
    Branch        & BEQ, BNE, BLT, BGE \\ \hline
    Jump          & JAL, JALR \\ \hline
    Upper imm.    & LUI \\ \hline
  \end{tabular}
\end{table}

\subsubsection*{Design decisions}
The following decisions were fixed in a shared interface specification before implementation
started, so that the five team members could develop their modules in parallel:
\begin{itemize}
  \item \textbf{Harvard memory organisation:} separate 4\,KB (1024-word) instruction and data
        memories with combinational read, which removes the structural hazard of a single memory port.
  \item \textbf{Branch resolution in EX} with a predict-not-taken policy; a taken branch or jump
        flushes the IF/ID and ID/EX registers, giving a two-cycle penalty.
  \item \textbf{Operand forwarding} from the EX/MEM and MEM/WB registers to both ALU inputs, with
        priority to the more recent EX/MEM value.
  \item \textbf{Load-use stall} of one cycle: the PC and IF/ID register hold their values and a
        bubble is inserted into ID/EX.
  \item \textbf{Register-file write-first bypass}, so that a value written in WB is read by the
        instruction in ID in the same cycle.
  \item \textbf{Branch flush has priority} over a load-use stall.
  \item \textbf{Jumps:} \texttt{JAL} and \texttt{JALR} write $PC+4$ to \texttt{rd}; the EX stage
        selects $PC+4$ instead of the ALU output for these instructions. \texttt{LUI} is executed as
        $0 + imm$ in the ALU.
\end{itemize}

\subsection{Instruction Fetch Stage}

The IF stage consists of the program counter register (\texttt{pc\_reg}), the instruction memory
(\texttt{instr\_mem}) and the PC selection multiplexer.

\textbf{Program counter.} \texttt{pc\_reg} is a 32-bit register that is cleared to address 0 on
reset. It has an enable input driven by the \texttt{pc\_write} signal of the hazard detection unit:
when a load-use hazard is detected, \texttt{pc\_write}~=~0 and the PC holds its value so that the
same instruction is fetched again in the next cycle.

\textbf{Next-PC selection.} The next PC is chosen by
\begin{equation}
  PC_{next} = \begin{cases}
     branch\_target, & \text{if } pc\_src = 1\\
     PC + 4,         & \text{otherwise}
  \end{cases}
\end{equation}
where \texttt{pc\_src} and \texttt{branch\_target} come from the branch unit in the EX stage.
Sequential fetch ($PC+4$) is therefore the default prediction.

\textbf{Instruction memory.} \texttt{instr\_mem} is a 1024-word array that is initialised with NOP
instructions (\texttt{addi x0, x0, 0}) and then loaded from a hexadecimal program file with
\texttt{\$readmemh}. Because all instructions are word-aligned, the word index is
\texttt{addr[11:2]} and the read is combinational, so an instruction is available in the same cycle
as its address.

\textbf{IF/ID register.} \texttt{if\_id\_reg} stores the fetched instruction and its PC. It has two
control inputs from the hazard detection unit: \texttt{en} (\texttt{if\_id\_write}) holds the
register during a stall, and \texttt{flush} (\texttt{if\_id\_flush}) replaces the instruction with
a NOP when a branch or jump is taken. Flush has priority over the stall, and reset also loads a NOP.

\subsection{Top-Level Integration}

The top module \texttt{riscv\_pipeline\_top} instantiates all stage modules, the four pipeline
registers, the forwarding unit and the hazard detection unit using the port names defined in the
shared interface specification. It also contains the two remaining multiplexers of the datapath:
the next-PC multiplexer described above and the write-back multiplexer
\begin{equation}
  wb\_data = \begin{cases}
     mem\_rdata, & \text{if } mem\_to\_reg = 1\\
     result,     & \text{otherwise.}
  \end{cases}
\end{equation}
The register source fields \texttt{rs1} = \texttt{instr[19:15]}, \texttt{rs2} = \texttt{instr[24:20]}
and the destination field \texttt{rd} = \texttt{instr[11:7]} are extracted from the IF/ID
instruction and sent both to the register file and to the hazard detection unit. The value
written back in WB is also used as the MEM/WB forwarding source, and the EX/MEM \texttt{result} is
used as the EX/MEM forwarding source and as the data-memory address.
